\section{Đạo đức và pháp lý}
\label{sec:dao-duc}

Đề yêu cầu tuân thủ điều khoản nguồn, \texttt{robots.txt}, giới hạn truy cập và quy
định dữ liệu cá nhân; dữ liệu có yếu tố con người cần văn bản đồng ý và ẩn danh
trước khi lưu \cite[tr.~2, mục II]{de}. Mục này ghi những gì nhóm đã làm và những
ranh giới nhóm tự đặt.

\subsection{Nguồn và quyền sử dụng}

\subsubsection{Corpus của A.2}

Bản \emph{Truyện Kiều} Trương Vĩnh Ký 1911 thuộc phạm vi công cộng theo thời hạn
bảo hộ. Nhóm chuyển sang bản này \emph{sau khi} rà soát cho thấy căn cứ quyền đối
với một bản số hóa khác (revision \texttt{206607}) chưa đủ --- việc công khai để
đọc không đồng nghĩa được phân phối lại. Kết luận ``chưa đủ căn cứ'' được giữ lại
trong tài liệu thay vì xóa, để người đọc thấy đường ra quyết định.

Trong quá trình audit kỹ thuật, nhóm phát hiện \texttt{robots.txt} chặn đường
\texttt{/w/} của nguồn đó. Sau khi phát hiện, nhóm \emph{không} lặp lại truy cập
API hay raw và chỉ dùng route \texttt{/wiki/} được cho phép. Giới hạn này được ghi
rõ, không che giấu.

\subsubsection{Nguồn thủ tục của Phần D}

Nhóm phân loại nội dung theo quyền trước khi đưa vào dữ liệu
(Mục~\ref{sec:d1}): chính văn quy phạm và văn bản hành chính không thuộc phạm vi
bảo hộ quyền tác giả, nên có căn cứ dùng làm nguồn soạn QA kèm dẫn chứng đầy đủ;
còn trang tin, bài hướng dẫn và video thì chỉ dùng để \emph{tìm} và đối chiếu văn
bản, không chuyển nguyên trạng vào tập train.

Về truy cập: tuân thủ \texttt{Crawl-delay: 5} của cổng Hà Nội, truy xuất tuần tự,
không crawl toàn cổng, không truy cập các đường dẫn hồ sơ cá nhân, nộp/sửa hồ sơ,
SSO hay đăng nhập mà \texttt{robots.txt} chặn. Khi một host trả HTTP 404 cho
\texttt{robots.txt}, nhóm ghi rõ rằng \emph{không có tệp robots không phải là giấy
phép}.

Có một cơ chế trong pháp luật Việt Nam cho phép dùng có điều kiện dữ liệu được bảo
hộ cho nghiên cứu và huấn luyện AI (khoản 5 Điều 7 Luật SHTT; Điều 37a--37c
NĐ 17/2023 sửa bởi 134/2026). Nhóm ghi nhận cơ chế này nhưng \textbf{không} dựa vào
nó để duyệt các hiện vật pilot: nó không phải giấy phép cào hay phân phối snapshot.

\subsection{Dữ liệu có yếu tố con người}

Hai nhóm dữ liệu người xuất hiện trong Phần~D: câu hỏi khảo sát cho tập test của
D.2, và nhãn ưu tiên của D.4.

\subsubsection{Quy trình đồng ý và ẩn danh}

\begin{itemize}[leftmargin=2em]
\item Thu dữ liệu người cần \textbf{văn bản đồng ý} trước khi thu.
\item Dữ liệu được \textbf{ẩn danh trước khi lưu} làm dữ liệu nghiên cứu.
\item \textbf{Phiếu đồng ý có định danh để riêng, không đưa vào Git.} Kho mã có
quy tắc bỏ qua \texttt{/consent/} và \texttt{/surveys\_raw/}.
\item Có mốc để người tham gia \emph{rút} dữ liệu.
\end{itemize}

\subsubsection{Phân quyền trong nhóm}

Nhóm tách quyền truy cập theo vai trò, không theo mức tin cậy:

\begin{itemize}[leftmargin=2em]
\item Cao Văn Hà phụ trách lưu trữ vận hành: kho riêng mã hóa và sổ token liên kết
định danh.
\item Hà Mạnh Thế và Trần Duy Hưng \textbf{không} nhận phiếu đồng ý, bảng liên kết
định danh, hay khảo sát thô. Họ chỉ nhận dữ liệu DPO đã được phân quyền, và
\textbf{không} xem tập \texttt{USER\_TEST} trước khi giao thức đánh giá yêu cầu.
\item Hai người này phải \emph{xác nhận riêng bằng văn bản} rằng họ chỉ nhận dữ
liệu DPO được phân quyền.
\end{itemize}

Lý do: người chấm nhãn không cần biết ai đã trả lời khảo sát để làm việc của mình.
Giảm phạm vi truy cập giảm rủi ro, và việc xem trước tập test sẽ phá phép đo ở D.5.

\textbf{Trạng thái \chualam:} chưa có người tham gia nào được mời, chưa thu phiếu
đồng ý nào, và chưa có nhãn ưu tiên nào được tạo. Bộ công cụ đã soạn, hoạt động
thật thì chưa diễn ra.

\subsection{Ranh giới nhóm tự đặt về AI tạo sinh}

Ngoài các ràng buộc của đề, nhóm tự đặt bốn ranh giới và ghi chúng thành quy tắc
nội bộ:

\begin{enumerate}[leftmargin=2em]
\item \textbf{Không dùng AI tạo sinh thay khảo sát người thật.} Câu hỏi test phải
đến từ người dùng thật, vì mục đích của tập test là đo khả năng tổng quát sang cách
người dân diễn đạt (Mục~\ref{sec:d2}).
\item \textbf{Không bịa} nguồn, QA đã kiểm chứng, khảo sát người thật, điểm người
chấm, kết quả huấn luyện hay mức đồng thuận.
\item \textbf{Nhận xét của AI không được tính vào quota kiểm thủ công.} Trong vòng
rà pilot, AI có đối chiếu QA với đoạn trích và đề xuất sửa từng câu, nhưng số QA
``được người duyệt'' vẫn giữ bằng $0$ cho tới khi người thật điền phiếu.
\item \textbf{Không điền vào cột kết luận của người kiểm.} Phiếu kiểm có cột mẫu
minh họa, nhưng cột kết luận của người kiểm để trống.
\end{enumerate}

Ranh giới thứ ba và thứ tư là lý do báo cáo này ghi ``0 QA được người duyệt'' dù đã
có 20 QA qua kiểm cấu trúc tự động.

\subsection{Những gì không đưa vào kho mã}

Theo quy tắc riêng tư và quy định của đề, kho mã \emph{không} chứa: secret và khóa
truy cập; dữ liệu cá nhân, khảo sát thô, phiếu đồng ý; checkpoint và mô hình nền.
Hệ quả là một số hiện vật pilot chỉ còn báo cáo JSON đã commit, và người tái lập
cần nhận bàn giao riêng (Mục~\ref{sec:tai-lap}).

Về năm bài báo trong \texttt{references/}: chúng đã được theo dõi trong Git từ
trước. Quyền phân phối lại chưa được kiểm chứng, nên nhóm ghi nhận đây là rủi ro và
xin quyết định trước khi công khai kho mã, thay vì tự xóa hoặc tự công khai.

\subsection{Giới hạn về tính chính xác của nội dung pháp lý}

Sản phẩm là trợ lý tra cứu, không phải nguồn pháp lý. Hai điểm được ghi vào đặc tả
đầu ra ngay từ đầu (Mục~\ref{sec:gioi-thieu}):

\begin{enumerate}[leftmargin=2em]
\item Câu trả lời phải \textbf{dẫn nguồn}, vì giá trị pháp lý thuộc về văn bản gốc.
\item Khi ngoài phạm vi hoặc chưa đủ căn cứ, trợ lý phải \textbf{nói rõ} thay vì
đoán.
\end{enumerate}

Hai mâu thuẫn tìm được ở Mục~\ref{sec:d1} --- thời hạn cấp trích lục giữa hai quyết
định, và việc trộn biểu phí chứng thực với biểu lệ phí hộ tịch --- cho thấy rủi ro
này không phải lý thuyết. Một hệ thống trả lời tự tin về thời hạn ``3 ngày'' mà
không nêu văn bản căn cứ sẽ gây hại thật cho người dùng.
